\documentclass[11pt,a4paper]{article}
\usepackage[T1]{fontenc}
\usepackage[utf8]{inputenc}
\usepackage{lmodern,microtype}
\usepackage[margin=25mm]{geometry}
\usepackage{amsmath,amssymb}
\usepackage{graphicx,booktabs,tabularx,array}
\usepackage{enumitem,caption,xspace,xcolor}
\usepackage[numbers,sort&compress]{natbib}
\usepackage[hyphens]{url}
\usepackage[colorlinks=true,linkcolor=black,citecolor=black,urlcolor=blue]{hyperref}
\usepackage{fancyhdr,placeins,listings}
\hypersetup{pdftitle={Missile Intercept: Modular Simulation Software with Explicit Data Contracts and Traceable Outputs},pdfauthor={Anonymous Authors}}
\definecolor{ink}{HTML}{142C44}
\definecolor{codebg}{HTML}{F4F6F8}
\captionsetup{font=small,labelfont=bf}
\setlength{\parskip}{3pt}
\setlength{\parindent}{1.2em}
\setlength{\emergencystretch}{2em}
\setlist{nosep,leftmargin=*}
\clubpenalty=10000
\widowpenalty=10000
\pagestyle{fancy}
\fancyhf{}
\fancyhead[L]{\small\textit{Missile Intercept: Software Architecture and Data Contracts}}
\fancyhead[R]{\small Anonymous preprint}
\fancyfoot[C]{\thepage}
\renewcommand{\headrulewidth}{0.3pt}
\setlength{\headheight}{14pt}
\lstset{basicstyle=\small\ttfamily,columns=fullflexible,keepspaces=true,breaklines=true,breakatwhitespace=true,showstringspaces=false,backgroundcolor=\color{codebg},frame=single,rulecolor=\color{black!15},framesep=6pt,aboveskip=10pt,belowskip=10pt}
\input{evidence_numbers}
\input{software_numbers}
\newcommand{\code}[1]{\texttt{#1}}
\newcolumntype{Y}{>{\raggedright\arraybackslash}X}
\newcolumntype{P}[1]{>{\raggedright\arraybackslash}p{#1}}
\title{\vspace{-1em}\textbf{Missile Intercept: Modular Simulation Software\\with Explicit Data Contracts and Traceable Outputs}}
\author{Anonymous Authors}
\date{November 13, 2026}

\begin{document}
\maketitle
\thispagestyle{plain}
\begin{abstract}
Research simulation software must connect heterogeneous components without losing the meaning, ownership, or history of the data they exchange. This paper describes the software architecture of \emph{Missile Intercept}, a Python simulation archive, with an emphasis on its domain-independent implementation mechanisms. The framework separates configuration, named state representation, execution coordination, module families, and output recording. We specify the metadata-based array schema, packing and unpacking behavior, bounded telemetry storage, data ownership, and export semantics, and relate these mechanisms to concrete source interfaces. An isolated evaluation selects the generic data-handling classes without importing simulation models. Across \SoftwareProbeCount synthetic interface probes, \SoftwareRequirementsMet targeted properties hold and \SoftwareLimitations stronger diagnostic requirements are not enforced. Twelve allocation cases agree exactly with the predicted float64 array payload. Existing provenance-tooling tests provide a separate record of \TestsPassed passed cases and \TestsSkipped skip. Four original diagrams explain the architecture, data representation, recording pipeline, and evaluation boundary. The resulting contribution is an implementation-centered software description and a reproducible study of its generic data infrastructure; historical simulation reports retain their unverified provenance status.
\end{abstract}
\noindent\textbf{Keywords:} scientific software; modular simulation; data contracts; telemetry; software evaluation

\section{Introduction}
Research simulation frameworks combine more than a numerical model. They must represent state consistently, connect modules with different responsibilities, record intermediate quantities, and preserve outputs in a form that another researcher can inspect. Errors at these boundaries can be difficult to distinguish from errors in the scientific model. An unexplained array offset, a mutable output view, or an undocumented export format can change how results are interpreted even when the underlying computation is unchanged.

\begin{figure}[t]
\centering
\includegraphics[width=\linewidth]{figures/architecture.pdf}
\caption{Software architecture of the archived framework. Solid connections show configuration or component dependencies and the coordinator's recording interface; dashed connections show output and presentation relationships. Named directories are concrete source locations. Historical report assets are shown as a presentation category, not as outputs proven to originate from the currently archived implementation.}
\label{fig:architecture}
\end{figure}

Software papers make these engineering choices assessable by describing interfaces and implementation alongside evidence. Stable-Baselines3 emphasizes a consistent user interface and implementation quality \citep{raffin2021}; CleanRL foregrounds implementation transparency and links reported results to inspectable experiments \citep{huang2022}. Benchmark papers add a further obligation: the tasks, evaluation protocol, and recorded runs must make comparisons interpretable, as illustrated by OpenML Benchmarking Suites \citep{bischl2021}. These examples motivate documenting what the present code actually provides before assigning it a benchmark role.

\emph{Missile Intercept} is an archived Python framework whose documentation describes a six-degree-of-freedom simulation with environmental, observation, decision, and command modules. Its original application is interception simulation. Here, we describe that application at the level of software organization and study the reusable data infrastructure in detail. The paper does not develop domain-specific decision laws or evaluate operational effectiveness.

The software contains mechanisms that warrant an explicit technical description: metadata-bearing state fields, schema-derived array slices, typed configuration objects, a central execution coordinator, an optional accelerated backend, a bounded telemetry recorder, and post-processing and provenance utilities. These mechanisms are scattered across source files and reports. Bringing them together exposes both their intended separation and their present limitations.

This paper contributes:
\begin{enumerate}
  \item A layered architecture and module-responsibility map connecting configuration, representation, execution, and output artifacts.
  \item An implementation-level account of the generic array schema and telemetry recorder, including validation order, complexity, ownership, and serialization behavior.
  \item A reproducible evaluation of isolated data-handling classes on dimensionless synthetic records, including diagnostic cases that do not satisfy stronger proposed invariants.
  \item A source-linked artifact package that distinguishes these new software observations from existing tooling tests and historical simulation pages.
\end{enumerate}
The contribution documents and examines the archived implementation. It does not claim a new integration algorithm, a standardized simulation API, or a new domain benchmark.

\section{Related Software and Positioning}
\label{sec:related}
The relevant comparison is with how scientific software exposes its design and supports its claims. Table~\ref{tab:position} summarizes four reference papers. Their application domains differ, so they are not performance baselines for this framework.

\begin{table}[t]
\caption{Related software and benchmark papers, and the corresponding reporting concern addressed here. The comparison concerns software description rather than relative performance.}
\label{tab:position}
\small
\begin{tabularx}{\linewidth}{@{}P{0.23\linewidth}YY@{}}
\toprule
Reference & Emphasis in the reference & Corresponding concern here \\
\midrule
Stable-Baselines3 \citep{raffin2021} & Consistent interfaces and reliable implementations & Identify concrete interfaces and separate executed tests from untested capabilities. \\
CleanRL \citep{huang2022} & Accessible implementation details and experimental records & Explain data-handling mechanisms and retain inspectable probe outputs. \\
OpenML suites \citep{bischl2021} & Explicit tasks, evaluation procedures, and reusable run records & State the evaluation unit and avoid treating unrelated historical pages as one benchmark. \\
JSBSim \citep{berndt2026} & A standalone flight-dynamics library and a statement of software need & Distinguish software architecture, module scope, and evidence of physical fidelity. \\
\bottomrule
\end{tabularx}
\end{table}

JSBSim provides a domain-relevant example of presenting a flight-dynamics software library through its purpose, architecture, and software ecosystem \citep{berndt2026}. Our archive has a narrower research-code scope: source-level module boundaries are visible, but a stable external plugin interface or broad compatibility record is not established. We therefore describe existing Python classes and files directly rather than presenting an unsupported interoperability claim.

The distinction between a software framework and a benchmark is also material. A benchmark needs more than a batch-execution script: it needs identified tasks, a fixed evaluation protocol, attributable runs, and an interpretation of the reported measurements. In the present archive, historical report pages lack matching execution-provenance records. This paper consequently adopts the software-paper form and evaluates generic software properties with a newly recorded protocol.

Project organization and scriptable research workflows are consistent with established scientific-computing practice \citep{wilson2017,sandve2013}. We use provenance as supporting infrastructure for the software description. It is not the principal scientific object of the paper.

\section{Design Requirements and System Architecture}
\label{sec:architecture}
\subsection{Design requirements}
The archived implementation addresses four practical software requirements.

\textbf{R1: explicit representation.} A numeric record should retain an interpretable relationship to field names, shapes, units, coordinate-frame labels, and ownership. Code should obtain slices from one schema rather than distribute unexplained offsets across modules.

\textbf{R2: separation of responsibilities.} Configuration, data representation, execution coordination, domain modules, and presentation should be distinguishable. This makes it possible to inspect a data-handling component without loading the full simulation stack.

\textbf{R3: bounded recording.} The recorder should have an explicit storage capacity, identifiable occupied rows, and visible behavior when input or storage constraints are violated. The caller must be able to distinguish an accepted row from a dropped row.

\textbf{R4: inspectable outputs.} Named columns and documented export behavior should support post-processing, while source and environment records should identify the scope of newly observed checks. These requirements concern software behavior; satisfying them does not establish the validity of a scientific model.

\subsection{Layered organization}
Figure~\ref{fig:architecture} provides the system overview. Configuration and schema definitions supply the representation boundary. The coordinator in \code{main.py} connects the simulation modules and numerical backend. The four domain directories contain model, observation, decision, and command families. Telemetry and post-processing form an output boundary, while batch and analysis utilities provide additional experiment-facing organization. The diagram is a software-responsibility view, not an exact function-call trace.

Table~\ref{tab:modules} maps this view to code. A static inventory contains \SimulationPythonFiles Python files under the simulation directory, including utilities and tests. The file count describes the archive's extent, not its maturity or test coverage.

\begin{table}[t]
\caption{Responsibilities and concrete software boundaries. Domain-specific algorithms are represented only by their module families.}
\label{tab:modules}
\small
\begin{tabularx}{\linewidth}{@{}P{0.28\linewidth}YY@{}}
\toprule
Location & Responsibility & Interface or artifact \\
\midrule
\code{config.py} & Configuration and typed data containers & \code{Scenario}, \code{SimConfig}, state/observation/command objects \\
\code{state\_schema.py} & Named metadata and array layout & \code{StateField}, \code{StateSchema}, named slices \\
\code{main.py} & Execution coordination & \code{Simulation} and recording calls \\
\code{integrator.py}, \code{fastcore.py} & Numerical backend and optional accelerated code & Array-oriented backend boundary \\
\code{vehicle/}, \code{sensors/}, \code{guidance/}, \code{control/} & Domain component families & State, observation, and command containers \\
\code{telemetry.py} & Bounded named records & \code{send}, \code{get}, \code{as\_dict}, \code{save\_csv} \\
\code{visualization.py}, \code{analysis.py}, \code{batch.py} & Presentation and experiment utilities & Saved data, plots, summaries \\
Root \code{tools/} and \code{tests/} & Independent artifact tooling & Environment JSON, snapshots, report inventories \\
\bottomrule
\end{tabularx}
\end{table}

\subsection{Configuration and execution boundaries}
The implementation uses dataclasses for configuration and exchanged records. \code{Scenario} describes a case, while \code{SimConfig} carries execution, recording, and output choices. Dataclass annotations make these structures inspectable but do not themselves enforce types or validate every field at runtime. Several defaults also originate from module-level configuration constants. The current design therefore combines explicit objects with global configuration dependencies.

The coordinator and numerical backend are separate source components. The orchestration source describes discrete module updates outside the internal continuous-integration stages. This paper records that architectural division without specifying domain update policies or evaluating numerical behavior. The optional accelerated path is a further implementation boundary: the presence of a Python fallback does not, on its own, demonstrate equal outputs or equal exception behavior across backends.

\section{Implementation of Generic Data Interfaces}
\label{sec:implementation}
\subsection{Field descriptors and schema construction}
\code{StateField} is a frozen dataclass with five attributes: \code{key}, \code{shape}, \code{unit}, \code{frame}, and \code{owner}. \code{StateSchema} stores an ordered tuple of descriptors, rejects duplicate keys, computes field offsets, and exposes a read-only slice mapping through \code{MappingProxyType}. Frozen descriptors and a read-only mapping constrain parts of the representation; the schema object as a whole is not declared immutable.

For $k$ fields, let $d_i$ be the product of field $i$'s shape dimensions, with $d_i=1$ for a scalar. The packed dimension and starting offsets are
\begin{equation}
D=\sum_{i=1}^{k} d_i, \qquad o_i=\sum_{j<i} d_j.
\end{equation}
Each field occupies the slice $[o_i,o_i+d_i)$. These relations specify a storage layout only. The generic example in Figure~\ref{fig:contracts} uses three dimensionless values and a scalar, without representing a physical system.

\begin{figure}[t]
\centering
\includegraphics[width=\linewidth]{figures/data_contracts.pdf}
\caption{Representation boundary implemented by the generic schema classes. Ordered metadata defines named array slices. Packing produces a flat record; unpacking produces copied array fields. The example is dimensionless synthetic data used by the accompanying probe harness.}
\label{fig:contracts}
\end{figure}

The schema computes a layout hash from a sorted-key JSON serialization of the ordered metadata records. Changing a descriptor's key, shape, unit, frame, or owner can therefore change the recorded layout identifier. Field order remains significant because the descriptors are serialized as a sequence. This identifier describes metadata consistency; it is not a checksum of stored values or a validation of their scientific meaning.

\subsection{Packing, unpacking, and ownership}
\code{pack(mapping)} first compares the supplied and expected key sets. It allocates a flat floating-point array, converts each supplied value to a NumPy array, checks the exact shape for nonscalar fields, and requires one value for scalar fields. It writes values into the schema-derived slices and rejects a packed record containing nonfinite entries. The implementation's scalar check is based on element count, so a one-element array can satisfy a scalar field.

\code{unpack(array)} checks the packed array shape and returns a dictionary. Nonscalar fields are reshaped and copied; scalar fields become Python floats. These are asymmetric contracts: the observed implementation checks finiteness during packing but not during unpacking. The distinction matters when external code supplies a packed array that did not come from a successful call to \code{pack}.

Listing~\ref{lst:schema} demonstrates the generic interface used in the isolated evaluation. The harness selects the class definitions from source and supplies their standard-library and NumPy dependencies; it does not import the application module.

\begin{lstlisting}[language=Python,caption={Dimensionless example using the extracted generic classes.},label={lst:schema}]
fields = [
    StateField("vector", (3,), "1", "synthetic", "demo"),
    StateField("label", (), "1", "-", "demo"),
]
schema = StateSchema(fields)
packed = schema.pack({"vector": [1., 2., 3.], "label": 7.})
record = schema.unpack(packed)
assert packed.shape == (4,)
\end{lstlisting}

For total packed dimension $D$, packing and unpacking have linear data-processing cost in $D$, in addition to key and field bookkeeping. Packing allocates a new array, and unpacking copies nonscalar fields. These costs follow from the implementation's array conversion and copying behavior; no throughput advantage is inferred.

\subsection{Bounded telemetry storage}
\label{sec:telemetry}
\code{TelemetryLogger} allocates a two-dimensional floating-point array with capacity $N$ and $C$ named columns, including time. It retains the occupied-row count $n$, the last accepted time, and an overflow counter. A successful \code{send(t, values)} writes one row and increments $n$. A width mismatch raises \code{ValueError}; nonfinite payload values raise \code{FloatingPointError} in the Python path; a full buffer increments the overflow counter and returns \code{False}. An equal or earlier timestamp also returns \code{False}.

Figure~\ref{fig:telemetry} separates validation, storage, and output ownership. Table~\ref{tab:api} gives the observable API. The caller should inspect the return value from \code{send}: a returned \code{False} is not an accepted record. The recorder drops rows at its fixed capacity rather than dynamically extending the data array.

\begin{figure}[t]
\centering
\includegraphics[width=\linewidth]{figures/telemetry_pipeline.pdf}
\caption{Generic telemetry implementation: record admission, fixed-capacity storage, access ownership, and CSV export. The two accessors have different aliasing behavior. The timestamp-finiteness boundary is an observed diagnostic limitation of the isolated Python path.}
\label{fig:telemetry}
\end{figure}

\begin{table}[t]
\caption{Generic public interfaces and their observed ownership or failure semantics.}
\label{tab:api}
\small
\begin{tabularx}{\linewidth}{@{}P{0.25\linewidth}YY@{}}
\toprule
Interface & Return / ownership & Validation or boundary \\
\midrule
\code{pack(values)} & New flat array & Exact keys and nonscalar shapes; scalar size; finite values \\
\code{unpack(packed)} & Dictionary; array fields copied & Packed shape checked; finiteness not checked \\
\code{send(t, values)} & Boolean acceptance; stores a row & Width, payload finiteness, capacity, and ordering checks in Python path \\
\code{get(name)} & View of an occupied column & Unknown name raises \code{KeyError}; returned view shares storage \\
\code{as\_dict()} & Copied occupied columns & Extra output memory proportional to occupied data \\
\code{save\_csv(path)} & Output path & Header and occupied rows; six decimal places \\
\bottomrule
\end{tabularx}
\end{table}

For the float64 array used in the experiment, its data payload is
\begin{equation}
B_{\mathrm{array}}=8NC\ \mathrm{bytes}.
\label{eq:memory}
\end{equation}
This is the allocated NumPy payload, not process memory. It excludes Python objects, allocator overhead, and copies made during post-processing. Appending an accepted row costs $O(C)$ for value checking and storage. Column lookup uses a linear search over names before creating a slice; \code{as\_dict} copies the occupied data and costs $O(nC)$.

\subsection{Export precision and backend scope}
CSV output contains the column-name header and only the occupied rows. Every numeric entry is formatted with six digits after the decimal point. The synthetic value $0.123456789$ is therefore written as \code{0.123457}. This is a presentation-oriented decimal representation, not an exact floating-point checkpoint. Documentation and downstream analyses should distinguish stored array values from exported decimal values.

The logger delegates recording to an accelerated helper when its availability flag is set. Our execution selects only the Python path and does not load that helper. Accordingly, the behaviors and boundary observations reported here apply to that isolated path. They do not establish cross-backend equivalence.

\section{Workflow, Artifacts, and Maintainability}
\label{sec:workflow}
The software architecture separates case/configuration objects, computation, recording, and presentation. At the data boundary, an inspectable workflow identifies the schema, passes named records through a defined interface, selects the required ownership when accessing outputs, and records the export format. This description exposes dependencies that a directory listing alone would not show.

The archive additionally supplies independent environment and artifact tools. These collect distribution metadata, record selected file digests, and inventory local report dependencies without importing simulation models. Their role is to describe retained artifacts. They do not reconstruct an older execution environment or establish that a historical page was generated by the current source. Practical guidance on preserving computational steps and versions supports keeping that distinction explicit \citep{sandve2013,barba2021}.

The current package is organized as research scripts with local imports, module-level defaults, and an optional acceleration dependency. A universal plugin mechanism, a packaged stable API, and arbitrary-backend interchangeability are not established. The contribution of the module boundaries is inspectability and separation within this archive. Stronger interoperability claims would require an explicit public contract and compatibility evidence.

The manuscript package preserves the implementation descriptions separately from measured observations. It includes source hashes, generic class-selection records, software-probe outputs, allocation measurements, existing tooling records, figure-generation code, and a clean LaTeX source archive.

\section{Software Evaluation}
\label{sec:evaluation}
\subsection{Protocol and environments}
We evaluate the generic data infrastructure rather than the domain model. The harness parses two original source files and selects only the class definitions \code{StateField}, \code{StateSchema}, \code{TelemetryLogger}, and \code{SimLoggerDefaults}. It records their source-file digests and selected line ranges. Dependencies are supplied explicitly, with acceleration disabled and synthetic column names. No simulation module, physical model, trajectory, or domain decision algorithm is imported or executed.

This form of isolation provides a precise evaluation unit: source-selected generic classes under a controlled namespace. It does not exercise the original module-import path or full application integration. The new run uses CPython \SoftwarePython and NumPy \SoftwareNumpy on Windows/AMD64. The previously recorded independent tooling suite uses CPython \AuditPython and pytest \AuditPytest. These are separate execution records, not a combined environment or a compatibility matrix.

The interface probes include valid records, malformed keys and shapes, nonfinite values, ownership checks, timestamp ordering, capacity overflow, and decimal export. Two diagnostic requirements deliberately ask for stronger finite-input enforcement at entry points where the implementation may not provide it. Allocation cases vary $N$ over four capacities and $C$ over three column counts, with no timing metric or process-memory measurement. All inputs are dimensionless and synthetic.

\subsection{Interface properties and diagnostic outcomes}
Of \SoftwareProbeCount recorded properties, \SoftwareRequirementsMet satisfy the targeted requirement. The remaining \SoftwareLimitations are informative boundary observations, not discarded cases. First, \code{unpack} accepts a correctly shaped array containing a NaN. Second, the isolated Python recorder accepts a NaN timestamp with finite payload values. A comparison-based ordering check does not itself enforce finite time.

These observations refine the interface description: successful packing enforces a stronger finite-value condition than arbitrary unpacking, and telemetry payload validation is narrower than validation of the entire record. Neither observation changes the original code. They should not be described as test failures against an established formal specification; they are unmet stronger requirements proposed by this study's diagnostic protocol.

Table~\ref{tab:probe-groups} groups the observations. The ownership checks are particularly relevant for consumers: modifying an array returned by \code{unpack} or \code{as\_dict} leaves the source storage unchanged, whereas modifying a column returned by \code{get} changes the underlying recorded array. Naming both behaviors avoids an unsupported assumption that all accessors return snapshots.

\begin{table}[t]
\caption{New synthetic interface evaluation. Detailed case names, observations, and extraction metadata are retained in \protect\path{evidence/software_probes.json}.}
\label{tab:probe-groups}
\small
\begin{tabularx}{\linewidth}{@{}P{0.25\linewidth}YY@{}}
\toprule
Probe group & Observed behavior & Interpretation \\
\midrule
Schema validation & Missing/extra keys, shape errors, duplicate names, and nonfinite packing rejected & Supports the stated packing contract \\
Schema ownership & Unpacked arrays are copied; metadata changes alter layout hash & Layout and aliasing behavior are inspectable \\
Telemetry admission & Width/nonfinite payload errors rejected; ordering and capacity drops observable & Return values and overflow counter matter \\
Access and export & Views differ from copies; CSV rounds to six decimal places & Consumers must choose ownership and precision explicitly \\
Stronger diagnostic conditions & Nonfinite unpacked input and NaN timestamps accepted & Two enforcement boundaries remain outside the current guards \\
\bottomrule
\end{tabularx}
\end{table}

\subsection{Array payload scaling}
All \AllocationCases allocation cases match Equation~\ref{eq:memory} exactly. Table~\ref{tab:memory} reports measured \code{data.nbytes} in KiB. Increasing capacity or column count increases the preallocated payload proportionally, as expected from the implementation's storage model. This experiment verifies the claimed array allocation, rather than demonstrating a speedup, a bound on total memory, or a scale limit for the full application.

\begin{table}[t]
\caption{Measured generic recorder array payload in KiB ($1\,\mathrm{KiB}=1024$ bytes), including the time column. Every entry equals $8NC/1024$ in the recorded environment.}
\label{tab:memory}
\centering\small
\begin{tabular}{@{}rrrr@{}}
\toprule
Capacity $N$ & $C=8$ & $C=32$ & $C=128$ \\
\midrule
256 & 16 & 64 & 256 \\
1024 & 64 & 256 & 1024 \\
4096 & 256 & 1024 & 4096 \\
16384 & 1024 & 4096 & 16384 \\
\bottomrule
\end{tabular}
\end{table}

\subsection{Supporting test and report records}
The original generic provenance-tooling suite has \TestsPassed passed cases, \TestsSkipped skipped case, and no recorded failure. The skip concerns directory symlink creation, which was unavailable in that environment. Ten earlier synthetic integrity scenarios match their expectations. These records remain separate from the new interface probes (Figure~\ref{fig:evaluation}); a skipped pytest case is not counted as a successful execution \citep{pytestskip}.

\begin{figure}[t]
\centering
\includegraphics[width=\linewidth]{figures/software_evaluation.pdf}
\caption{Evidence from the isolated data-interface evaluation and the separately recorded original tooling suite. Diagnostic gaps and platform skips retain distinct meanings. Neither group is a measurement of simulation performance.}
\label{fig:evaluation}
\end{figure}

Static inspection identifies \SimulationTestDefinitions simulation test-function definitions across \SimulationTestFiles files, but those tests are not executed in this study. The three historical HTML reports contain \ReportAssets inventoried local assets and have no matching report-specific provenance manifests. They therefore remain \code{legacy\_unverified}. Their numerical results are excluded from the newly measured software evaluation.
\FloatBarrier

\section{Discussion and Limitations}
\label{sec:limitations}
The software description exposes useful tradeoffs. Metadata and named slices make representation explicit, while contiguous arrays support numerical interfaces. Fixed-capacity recording makes the primary payload predictable, while copy-producing accessors add memory beyond that payload. Python fallbacks make some generic components inspectable in isolation, while accelerated paths remain separate verification targets. These are consequences of the existing implementation, rather than evidence that one design dominates alternative libraries.

The synthetic study has deliberately limited external validity. Class extraction preserves the selected definitions but replaces their surrounding application namespace; it does not test import-time behavior or integration with the application. The cases cover explicit interface properties, not all values, shapes, exceptions, or combinations. Two stronger diagnostic conditions are unfulfilled. Allocation measurements cover the array payload only. We report neither runtime throughput nor total process memory, and no comparative performance baseline is established.

The framework's architecture is described from a local archive without Git metadata. An inventory of \SelectedFiles selected files identifies the studied bytes but is not a full archive image or an authenticated execution history. Existing CI configuration specifies intended generic-tooling jobs, not completed cross-platform results. Scientific models, original simulation tests, accelerated-path equivalence, and the historical report-generation chain remain outside the executed study.

These limits determine the paper's claims. The archive supports an implementation-centered software description and a reproducible study of generic interfaces. It does not currently support treating historical pages as a standardized benchmark, attributing performance gains to a particular method, or asserting validated physical or operational capability. The module and data-interface account remains useful without making those stronger claims.

\section{Conclusion}
We presented \emph{Missile Intercept} through its software architecture, explicit array contracts, bounded recording mechanism, and traceable output interfaces. Four diagrams connect the source structure to representation and recording behavior. An isolated synthetic study records \SoftwareRequirementsMet satisfied properties and \SoftwareLimitations unenforced diagnostic conditions across \SoftwareProbeCount probes; \AllocationCases array allocations match the stated payload model. The package provides the code-level explanations and machine-readable observations needed to assess those claims, while preserving the distinction between generic software evidence and historical simulation results.

\section*{Software and Artifact Availability}
The accompanying folder contains the anonymous manuscript, BibTeX references, vector figures, all probe and build scripts, machine-readable results, and unchanged generic provenance tools with their original MIT license. The original archive is needed to repeat source extraction; rendering and compiling the delivered paper use the included artifacts. Anonymous authorship is retained during preparation. Author identities, affiliations, funding, and experimental history are not inferred from the archive.

\bibliographystyle{unsrtnat}
\bibliography{references}

\appendix
\section{Reproducing the Generic Software Study}
\label{app:reproduce}
The commands below assume that the original archive is the parent of the paper directory. The extraction command accepts \path{--source-root} for another location. Use a working copy when regenerating evidence so the delivered observations remain preserved.

\begin{lstlisting}
python scripts/probe_data_contracts.py --source-root ..
python scripts/draw_figures.py
python scripts/build_paper.py --verify-zip
\end{lstlisting}
The new synthetic study requires NumPy. Figure generation requires ReportLab; compilation requires pdfLaTeX and BibTeX. The data-interface harness records the source hashes, selected class ranges, environment, and case-level outcomes in \path{evidence/software_probes.json}. It does not import the simulation package. The source ZIP contains only the referenced manuscript inputs and four figure PDFs.

The existing generic-tooling result can be repeated separately:
\begin{lstlisting}
cd reproducibility
python -m pytest -q tests -p no:cacheprovider --junitxml=../evidence/tooling_tests.xml
cd ..
\end{lstlisting}
The pytest invocation is one command. Its outcome can differ when platform capabilities or dependencies differ; new execution records should not overwrite the interpretation of an older environment.

\section{Source and Evidence Map}
\begin{table}[h]
\caption{Implementation claims and the accompanying records. Original source paths are relative to the source archive; evidence paths are relative to this paper folder.}
\small
\begin{tabularx}{\linewidth}{@{}P{0.36\linewidth}Y@{}}
\toprule
Description or observation & Source or retained record \\
\midrule
Schema descriptors, slices, packing and unpacking & \path{missile_intercept/state_schema.py}; extraction hashes in \path{evidence/software_probes.json} \\
Recorder admission, ownership, export & \path{missile_intercept/telemetry.py}; same extraction record \\
Coordinator and configuration names & \path{missile_intercept/main.py}, \path{missile_intercept/config.py} \\
New synthetic probes and allocations & \path{scripts/probe_data_contracts.py}, \path{evidence/software_probes.json} \\
Architecture and interface diagrams & \path{scripts/draw_figures.py}, \path{figures/} \\
Original source inventory & \path{evidence/source_inventory.json} \\
Existing generic-tooling observations & \path{evidence/tooling_tests.xml}, \path{evidence/integrity_probes.json} \\
Historical report provenance status & \path{evidence/report_audit.json} \\
Build and PDF checks & \path{evidence/build_check.json}, \path{evidence/pdf_check.json} \\
Delivery integrity & \path{package_manifest.json} \\
\bottomrule
\end{tabularx}
\end{table}
\end{document}
